\chapter{To the instructor}
\label{ch:instructor}

You are being asked to teach a class that does not yet appear in the ACM/IEEE
curriculum guidelines under that name.
This chapter is the briefing.

\section{The course in one paragraph}

Students learn to treat a \emph{synthesized gate netlist} as the universal
intermediate representation of reconfigurable computing, to lower that netlist
onto a \emph{tensor engine} (here: Apple Silicon \texttt{MPSGraph}), and to
evaluate it either as a boolean-faithful oracle or as torus FHE with
certificates.
They then learn a second compiler direction---melting discrete hardware into
continuous tensors---and a third: designing ciphers that are born under
adversarial pressure.
The moral of the course is not ``FHE is fast on a laptop.''
The moral is: \emph{reconfigurable computing has a new fabric, and honesty about
what that fabric proves is part of the engineering.}

\section{Prerequisites}

\begin{center}
\begin{tabular}{@{}lp{0.62\textwidth}@{}}
\toprule
\textbf{Required} & Digital logic (LUTs, flip-flops, timing)~\cite{harris2012digital};
linear algebra; a first crypto course (modular arithmetic, attack models)~\cite{paar2009crypto}. \\
\textbf{Helpful} & One HDL (Verilog or VHDL); one systems language; Yosys as a
black box~\cite{yosys}. \\
\textbf{Not required} & Prior FHE, prior Metal, prior Swift, prior lattice
cryptography. The torus chapter is self-contained at the level this course needs.
\\
\bottomrule
\end{tabular}
\end{center}

\section{Two tracks}

Not every department can put an M-series Mac in front of every student.

\begin{description}
  \item[Systems track.]
        Apple Silicon laboratory.
        Students build, tick, and certificate-read HELUT from
        \corpus{REPRODUCE.md}.
        Labs~\ref{ch:lab} assume this track.
  \item[Theory track.]
        No GPU required.
        Students work the five-cell test, noise lemmas, and published logs as
        \emph{data}.
        A department can still grade a claim-audit final (below).
\end{description}

A mixed offering works: theory problem sets for everyone; systems labs as a
project option or a second section.

\section{Suggested grading}

\begin{center}
\begin{tabular}{@{}lp{0.22\textwidth}p{0.48\textwidth}@{}}
\toprule
Weight & Instrument & Why it exists \\
\midrule
30\% & Problem sets & Negacyclic arithmetic, packing, certificate reading. \\
40\% & Laboratories or theory equivalents & Receipts, not vibes. \\
20\% & \textbf{Claim audit} & Take a public sentence; run the five-cell test
from Chapter~\ref{ch:living}. Fail the assignment if the student upgrades a
hedge to a theorem. \\
10\% & Frontier essay & Pick one avenue from Chapter~\ref{ch:frontier}.
Must label it frontier. Must propose the \emph{first experiment} that would
graduate it. \\
\bottomrule
\end{tabular}
\end{center}

The claim audit is the signature assignment of this course.
It trains the skill the field currently lacks: refusing to launder an
open hedge into a keynote.

\section{What to lecture, what to assign as reading}

\begin{itemize}
  \item \textbf{Lecture} Parts~II--III (foundations and Pillar~I), the
        honesty chapter, and Theorem~\ref{thm:tensorlut} (\cid{19}).
  \item \textbf{Seminar} Pillar~II empirics (shatter / hold, campaign as
        case study) and Pillar~III in edition~\corpusedition{}.
        Do not lecture shatter grades as if they were Theorem~\ref{thm:tensorlut}.
  \item \textbf{Do not lecture as fact} Chapter~\ref{ch:frontier}.
        Give it as the last two seminars, with the frontier boxes left on
        the slides.
\end{itemize}

\section{How to pull a new edition mid-semester}

Science will move under your feet.
That is the point of a living textbook.

\begin{enumerate}
  \item Diff \corpus{directives/claim-sheet.md} against the epoch printed on
        the title page.
  \item If a new \cid{} row appeared, assign the matching lab receipt as a
        one-page addendum; do not pretend the old PDF contained it.
  \item If an \hid{} closed, tell the class explicitly: ``this asterisk is
        gone; here is the log.''
  \item If only frontier text moved, do not change the midterm.
\end{enumerate}

Bump \texttt{\textbackslash livingepoch} in \corpus{textbook/preamble.tex}
when you accept an addendum.
Cite git.

\section{Classroom norms}

\begin{itemize}
  \item Trivial / mock torus encodings are a \emph{shape laboratory}, not FHE.
        Students who call the boolean oracle ``encrypted computing'' lose
        points on purpose.
  \item \textbf{Campaign cryptanalysis} (Enigma / P1030680) is a \emph{case study} in
        Pillar~II methods, not the FHE claim, and not a decrypt.
  \item \textbf{Do not lecture} HELUT as an unsupervised or mobile agent, a
        self-rewriting FPGA, or a production encrypted CPU. Those sentences fail
        the five-cell test of Chapter~\ref{ch:living}.
  \item \hid{4} Grade~B is covering Track~A (\cid{52}--\cid{54}), not native
        $k{=}1$, not \texttt{cryptoPublicMS}, not PicoRV covering at production
        $N$.
  \item Speculative systems (encrypted LLM firewalls, self-modifying dark
        binaries, honey-token ledgers) are ethics-and-design discussions,
        not implementation homework, until they have \cid{} rows.
\end{itemize}

\growswhen{A published instructor's guide with slide skeletons, a Moodle/Canvas
quiz bank, and a non-Apple reference implementation of the negacyclic oracle
for the theory track.}
